\section{Appendix D. Reproduction and Versioned Materials}
\phantomsection
\label{app:d}
\label{sec:extension-provenance}
\label{sec:display-archive-r14}

The public \href{https://github.com/AlfredJamesLi/chinese-skillspan-benchmark/blob/main/REPRODUCIBILITY.md}{reproduction guide} separates three tasks: inspecting released results, scoring saved predictions, and rerunning training or inference. The \href{https://github.com/AlfredJamesLi/chinese-skillspan-benchmark/blob/main/reproduction/PAPER_INDEX.md}{paper-to-file index} links each table and figure to its supporting materials. The \href{https://github.com/AlfredJamesLi/chinese-skillspan-benchmark/blob/main/PAPER_NAMES.md}{resource naming guide} maps manuscript terms to exact file names and model identifiers.

Shared-guideline inference, Qwen adaptation, and the JobBERT supervision comparison retain their own data splits, protocol versions, and selection records. Supplementary tables and experimental notes provide detailed results and reproduction requirements, with Qwen adapters and frozen predictions available in the \href{https://doi.org/10.5281/zenodo.22851581}{model-reproduction archive}. Its model index links each adapter to its selected checkpoint, protocol, and recorded result. The annotation-quality documentation reports independently recomputed agreement statistics and the QA150 nominal-alpha correction; the original recent coder exports are retained by the authors. The \href{https://github.com/AlfredJamesLi/chinese-skillspan-benchmark/tree/main/reproduction/benchmark_profile}{Figure~\ref{fig:length-distributions-r40} reproduction files} provide the measured length-frequency counts, type counts, plotting script, and source-file hashes without recruitment wording.

Expanded-Silver results are pinned to \href{https://github.com/AlfredJamesLi/chinese-skillspan-benchmark/commit/11355ff567d647cb087b0f033e09fd0825dac3c4}{the September 17 result snapshot}. Its Table H maps to Table~\ref{tab:expanded-silver-main} (pooled runs) and Table~\ref{tab:expanded-silver-subsets} (subset runs). The Codex-pool JSON snapshots provide full-precision Gold150 scores and split counts. The model archive supplies the adapters, split identifiers, inference prompts, parser, and scoring commands. Retraining additionally requires the corresponding complete Silver training texts and labels, which are not included in that archive. The shared-guideline scoring protocol is separate from the historical JSON-offset protocol.

\subsection{Collection and Preparation Records}
The available public-institution records retain announcement titles, source names, URLs, listed publication dates, and geographic metadata where available. Listed publication dates are not collection dates. A complete source-by-stage census of advertisement counts, crawl windows, exclusions, and sentence-to-document mappings has not been recovered.

\subsection{Audit of Frozen Expanded-Silver Splits}
A post hoc audit of the complete frozen expanded-Silver files found that their hashes, counts, and identifier order match the archived split manifests. No identifiers or exact/NFC-normalized sentences are shared across partitions or with the 150-sentence human reference. An additional comparison using NFKC normalization, whitespace removal, and case folding identified the same single training--validation pair in both variants, differing only in full-width versus half-width semicolons; both records have empty span annotations. Three additional punctuation-variant pairs occur within each training set. No corresponding normalized-text overlap was found between training or validation and the Silver test partitions, or between any partition and the human reference. The frozen splits and reported results were retained; the audit does not establish that the validation overlap had no effect. These checks do not establish document-level or semantic near-duplicate isolation.

\subsection{Archive and Checkpoint Versions}
The manuscript cites \href{https://doi.org/10.5281/zenodo.22698504}{Zenodo v0.1.3}. Earlier snapshots remain at \href{https://doi.org/10.5281/zenodo.22685143}{v0.1.2} and \href{https://doi.org/10.5281/zenodo.22288338}{v0.1.1}; they do not contain the human reference or the released B2 training/development files. The corresponding GitHub tags v0.1.0 and v0.1.1 have been withdrawn. The first Zenodo snapshot remains historical and is not the recommended reproduction entry point. The released Hugging Face model-card revisions are \texttt{b40e930} for JobBERT-zh initialization and \texttt{5e68d49} for JobBERT-zh B2 continuation. The current model cards update documentation; the cited revisions identify the earlier release descriptions.

\subsection{Resource Access and Sampling Scope}
\label{sec:resource-access}
\input{tex/review20260923/access_matrix}
The September 21 audit verified the complete expanded files against their recorded hashes. The files remain author-held; these checks establish file identity, but neither public access nor successful retraining.


\paragraph{Expansion sampling.}
Preparation retained 2,064 valid unique sentences from the original supervision files and selected 7,936 new candidates: 3,174 cloud, 2,381 public-sector, and 2,381 listed-company sentences. Selection prioritized requirement-related keywords within each source and used fixed hash ranking with seed 20260913. Processing removed HTML and whitespace, split advertisements into sentence units, and kept normalized lengths of 8--400 characters. Matching sentences were excluded against existing corpus splits, supervision, human-review sets, and the available 3.2-million-sentence dump used for the paper's domain pretraining. The comparison key applies NFC, removes whitespace, and trims specified edge punctuation. Annotations use the prepared strings. These checks concern sentence matches under that key and do not establish advertisement-level or semantic near-duplicate isolation.

The source manifest and executable comparison-key function document these operations. The audit records 3,200,000 lines in the paper-main domain-pretraining dump, but its content hash was not retained in the audit. Other local pretraining dumps were checked but not used as exclusions: the selected set has 143 matches in a listed-company mixture and 4,919 in a domain mixture. These are distinct from the paper-main pretraining dump. The records do not establish exclusion from all possible pretraining sources.

\input{tex/review20260923/source_stages}

\subsection{Annotation Channels and Frozen Targets}
The original Silver-plus generation used Handbook B and its stage-specific prompt. The remaining 2,351 original LLM-annotated records used the v4.2.12-based prompt. Amendments in v4.2.13--v4.2.14 are stored separately; Table~\ref{tab:execution-version-matrix} gives the full version map. For expansion, Codex CLI and a third-party API-compatible proxy served as distinct annotation channels through the public annotation prompt. The final variant replaced proxy annotations for 2,500 candidate sentences with Codex annotations. Unresolved annotations were excluded, producing pools of 9,646 and 9,540 records. These names identify access channels rather than independently verified underlying models. The original Qwen split records also give sampling origins and type distributions for training and development examples.

Each experiment retains its executed protocol and original labels. The archived rule-case check identifies the documented experience-boundary example in the Silver layer rather than Gold150; it is a single-case check, not a corpus-wide compatibility result. Applying revised rules requires a separately versioned annotation release. The frozen shared-guideline inference prompt remains distinct from the original bulk-annotation prompt and later public Silver prompt.

\paragraph{Historical implementation.}
For some JobBERT runs, the exact tokenizer revision, special-token handling, truncation or windowing counts, and affected reference-span counts have not been recovered. The maximum input of 256 tokens is not a character limit, and input-length distributions alone cannot quantify truncation. The reported reference scores retain the full target sets. Per-seed checkpoint-selection records are stored with the available implementation history, including the inherited checkpoint's unresolved selection history.
